\chapter{Introdução}

O estudo da resposta de circuitos lineares no domínio da frequência é um pilar fundamental para a compreensão e o projeto de sistemas de processamento de sinais, permitindo a seleção de faixas de operação específicas e a atenuação de componentes indesejadas \cite{nilsson}. Dentre as diversas topologias de filtragem, os filtros ativos destacam-se por utilizarem amplificadores operacionais em conjunto com redes de resistores e capacitores. Essa configuração possibilita a síntese de polos complexos conjugados sem a necessidade de indutores (que costumam ser volumosos e não ideais em baixas frequências), além de fornecer isolamento entre estágios e ganho de tensão ao sinal.

Nesta quinta prática de laboratório, realiza-se o projeto integral, a análise teórica e a validação experimental de um filtro ativo passa-baixa utilizando a aproximação de Butterworth de 4ª ordem. A escolha da resposta de Butterworth justifica-se pela sua característica de banda passante maximamente plana, que garante uma distorção mínima de amplitude na região de interesse \cite{nilsson}. O projeto estabelecido possui requisitos estritos de projeto: uma frequência de corte ($f_c$) de $50\text{ Hz}$ e um ganho em banda passante ($K$) igual a 4.

O trabalho foi estruturado em etapas metódicas e complementares. Inicialmente, o projeto parte da determinação do filtro-protótipo normalizado. Em seguida, aplicam-se as técnicas de escalonamento de frequência e de magnitude para a obtenção dos valores analíticos dos componentes passivos, os quais são posteriormente ajustados para os valores comerciais mais próximos disponíveis em bancada. Toda a análise nodal do circuito no domínio de Laplace e a obtenção das expressões para os diagramas de Bode foram processadas com o auxílio computacional do software de álgebra simbólica Maxima \cite{maxima}.

A validação teórica pré-bancada pôde ainda ser suportada pela simulação do esquemático elétrico no ambiente SPICE \cite{ltspice}. Por fim, a montagem física na matriz de contatos permitiu a excitação do sistema com um gerador de sinais e a medição experimental da magnitude da função de transferência utilizando o osciloscópio, estabelecendo o comparativo rigoroso entre o modelo matemático e a resposta real do filtro projetado.

\section{Objetivos}

\begin{itemize}
	\item Projetar um filtro-protótipo ativo passa-baixa Butterworth de 4ª ordem e aplicar os devidos procedimentos de escalonamento de frequência e de magnitude para atingir $f_c = 50\text{ Hz}$.
	\item Ajustar os componentes analíticos para valores comerciais, recalculando teoricamente a função de transferência, a nova frequência de corte e o ganho real do filtro.
	\item Aplicar a transformada de Laplace na análise nodal de circuitos lineares com amplificadores operacionais \cite{nilsson}.
	\item Praticar o uso de softwares de computação algébrica e de simulação em circuitos eletrônicos para a resolução de sistemas de equações algébricas e levantamento do gráfico de Bode \cite{maxima}.
	\item Desenvolver competências práticas de instrumentação em laboratório, utilizando o gerador de funções e o osciloscópio para o rastreamento em frequência do circuito físico, validando as características de atenuação e ganho previstas pelo projeto.
\end{itemize}